% ---------
% --------
%!TEX TS-program = pdflatex
%!TEX encoding = UTF-8 Unicode
%!TEX spellcheck = English (Aspell)

\documentclass[11pt]{article}
\usepackage{lucy,handout,graphicx}
\usepackage{fourier-orns}
\usepackage{mathtools}
\usepackage[normalem]{ulem} % either use this (simple) or
\usepackage{tikz}

\def\Cdot{\mathbin{\text{\normalfont \textbullet}}}
\def\Sym#1{\texttt{\color{BrickRed}#1}}
\def\BlueSym#1{\textbf{\texttt{\color{RoyalBlue}#1}}}
\allowdisplaybreaks

\begin{document}

\begin{center}
\Large\textbf{CSCI 332, Fall 2026}%
\\
\LARGE\textbf{Homework 1}%
\\[0.5ex]
\large Due Monday, August 31 Anywhere on Earth (6am Tuesday)
\end{center}

\bigskip
\hrule
\bigskip

\subsection*{Submission Requirements}
\begin{itemize}
    \item Type or clearly hand-write your solutions into a PDF format so that they are legible and professional. Submit your PDF on Gradescope. 
    \item Do not submit your first draft. Type or clearly re-write your solutions for your final submission. If your submission is not legible, we will ask you to resubmit.
    \item Use Gradescope to assign problems to the correct page(s) in your solution. If you do not do this correctly, we will ask you to resubmit.
\end{itemize}

\subsection*{Academic Integrity}

Remember, you may access \EMPH{any} resource in preparing your solution to the homework. However, you \EMPH{must}
\begin{itemize}
    \item write your solutions in your own words, and
    \item credit every resource you use (for example: ``Bob Smith helped me on
    this problem. He took this course at UM in Fall 2020''; ``I found a solution
    to a problem similar to this one in the lecture notes for a different
    course, found at this link: www.profzeno.com/agreatclass/lecture10''; ``Here is a link to the full transcript of my chat with ChatGPT about the problems on this HW''.)
    Ask if you need any clarifications about how to properly cite your sources!
\end{itemize}




\newpage
%----------------------------------------------------------------------

\headers{CSCI 332}{Homework 1 (due August 31)}{Fall 2026}


\begin{enumerate}
%\parindent 1.5em \itemsep 4ex plus 0.5fil

%----------------------------------------------------------------------

\item (0 points) List the resources you used for each problem in this homework. If you used no resources except
lecture notes and videos and the linked textbook, you can say ``none''. If you used
an AI tool, you should put a link to your full conversation here, and you must use the provided prompt to guide the AI to act
as a tutor, not an answer-provider.

You will be asked to resubmit your homework if this section is blank.

Remember that no matter what resources you consulted, you must write your solutions
\emph{in your own words}.

\item 

For your upcoming family reunion, you are in charge of one of the day's activities.
You decide to have the family head out in two-person canoes. However, you know that
you need to be careful in pairing your family members up, because they can get jealous
if they get stuck spending time with one person when they would prefer to spend time with another.
Because you are taking CSCI 332, you wonder if you could re-work the stable matching
definition from the hospital-student problem to work for your canoe pair problem. You consider
asking each family member to order the rest of the family members according to their
preference for a canoe partner.

For example, for $n=4$ family reunion attendees, you might get the following list:

\begin{center}
\begin{tabular}{ c c c }
    Alice: Ben, Cleo, Desiree \\
    Ben: Cleo, Alice, Desiree \\
    Cleo: Alice, Ben, Desiree\\
    Desiree: Alice, Ben, Cleo
\end{tabular}
\end{center}


\begin{enumerate}[(a)]
    \item (3 points) Re-work the definition of stable matching from the hospital-medical student
        version of the problem to the canoe pair problem.
        You should describe the input of the problem and the desired output of the problem.
    \item (3 points) Does a stable matching always exist for the canoe pair problem?
    \item (3 points) How do you know?
\end{enumerate}

\item (1 point) Back to the hospital-student stable matching problem: describe
and analyze an efficient algorithm to determine whether a given set of preferences
has a \emph{unique} stable matching.

\end{enumerate}
%%%%%%%%%%%%%%%%%%%%
\end{document}
